
\subsubsection{2.1 Reinforcement Learning from Execution Feedback}

RLEF with binary execution rewards was established by Gehring et al. [2024], who demonstrated order-of-magnitude inference-time sample reductions through binary pass/fail rewards combined with GRPO policy optimization. That work confirmed that binary execution correctness is sufficient for meaningful policy improvement and established the MBPP training $\rightarrow$ HumanEval+ evaluation protocol used here.

The mathematical foundation for why binary rewards lead to gradient starvation is the GRPO standard deviation identity: $\sigma = \sqrt(k(G-k))/G$, where k is the number of correct completions out of G. When k=0 or k=G, $\sigma =0$ and no gradient is produced. This identity underlies both the motivation for variance-guided selection and the cold-start diagnosis.

\subsubsection{2.2 Gradient Starvation in GRPO}

The empirical consequences of gradient starvation were quantified by Nie et al. [2026]: at G=4 with binary rewards, 54.75\% of training groups have all-fail outcomes and 14.50\% have all-pass outcomes, producing 69.25\% zero-gradient groups in expectation. The experiments reported here provide a per-problem characterization of this effect for MBPP under constrained generation, finding a more severe all-fail rate (91.7\%), suggesting that generation parameter constraints amplify practical gradient starvation beyond theoretical estimates for this model.

\subsubsection{2.3 Online Data Selection for GRPO Efficiency}

Multiple recent methods converge on selecting intermediate-difficulty problems as the key to GRPO efficiency --- all as online methods operating during training.

Jiang et al. [2026] (VIGOR, arXiv:2607.22002) allocate rollouts based on variance-utility signals computed during active training. This method parallels offline profiling conceptually but requires training loop integration, precluding the decoupled ''profile once, train anywhere'' design.

Baroian and Berger [2026] (Prompt Replay, arXiv:2603.21177) track per-problem running pass rates during training and prioritize problems with pass rate approximately 0.5 --- the maximum-variance point under the binary variance identity. This is the closest conceptual predecessor to the method proposed here, but it cannot operate as a pure offline step.

Sun et al. [2025] demonstrate difficulty-targeted online selection for math reasoning with scalar rewards, achieving 23--62\% compute reduction. This work targets a different reward type and domain; the data selection intuition is extended here to binary-reward code generation with an offline design.

Cui et al. [2026] (LZE, arXiv:2605.17003) fuse pass-rate momentum and outcome-uncertainty during training. Like the methods above, this requires an active training loop.

The offline approach in this paper occupies a distinct niche: profile once with a frozen model, train anywhere. The cold-start failure identified here --- that parameter mismatch between profiling and training can nullify the entire selection --- is a failure mode that online methods naturally avoid, since they compute difficulty signals during training where profiling-training alignment is not a concern.

\subsubsection{2.4 Positioning}

This work contributes: (1) the first offline variance-guided selection method for binary-reward code generation RLEF, (2) to the best of our knowledge, the first systematic per-problem binary reward variance distribution characterization on MBPP (prior GRPO-on-MBPP papers encounter this distribution but do not report it at per-problem granularity), and (3) the first systematic identification of cold-start as a necessary precondition for data selection in RLEF.

